\chapter{Choosing a Data Structure}\label{ch:ds-basics}

A \defn{data structure} is a way of organising data so that the operations an algorithm needs are fast. Choosing the right
one is often the entire solution: the same idea runs in $\Oh{n^2}$ with a plain array and in $\Oh{n \log n}$ with a heap, and
the heap version is only three lines longer. The previous chapter showed how to use the STL containers; this chapter is about
the thinking behind them: which operations each structure makes cheap, which it makes expensive, and how simple structures
(stacks, queues, heaps, sets) combine into solutions for problems that look nothing like textbook exercises.

\begin{chapterintro}{\bronze\ \silver\ \gold}{Chapters~\ref{ch:complexity} and~\ref{ch:cpp}.}
  \item compare structures by the cost of each operation, and pick one from the operations you need;
  \item use stacks for nesting and ``nearest'' questions, including the monotonic stack;
  \item use queues and deques for arrival order and sliding windows;
  \item use heaps for repeated minimums, scheduling and running medians;
  \item maintain sets of values and sets of intervals, and simulate deletions with array-based linked lists;
  \item augment simple structures (a stack with its minimum, a queue made of two stacks).
\end{chapterintro}

% ======================================================================================================================
\section[Operations decide]{Operations decide \hfill \bronze}\label{sec:ops-decide}

Before choosing a structure, list the operations your algorithm performs, with how often it performs each: ``$n$ inserts,
$q$ queries for the smallest element, occasional deletions''. Then pick the structure that makes the \emph{frequent}
operations cheap. Figure~\ref{fig:ds-costs} compares five ways of storing a changing set of numbers.

\begin{figure}[htbp]
\centering
\newcommand{\cf}[1]{\cellcolor{fggreen!22}#1}
\newcommand{\cl}[1]{\cellcolor{fgorange!22}#1}
\newcommand{\cs}[1]{\cellcolor{fgverm!22}#1}
\newcommand{\cx}{\cellcolor{fgline!25}\textcolor{fggray}{---}}
\small
\setlength{\tabcolsep}{5pt}
\begin{tabular}{l*{6}{c}}
\toprule
\thd{Structure} & \thd{insert} & \thd{delete $x$} & \thd{find $x$} & \thd{minimum} & \thd{$k$-th smallest} & \thd{sorted scan} \\
\midrule
unsorted array & \cf{$1$} & \cs{$n$} & \cs{$n$} & \cs{$n$} & \cs{$n$} & \cl{$n \log n$} \\
sorted array & \cs{$n$} & \cs{$n$} & \cl{$\log n$} & \cf{$1$} & \cf{$1$} & \cf{$n$} \\
balanced tree (\texttt{set}) & \cl{$\log n$} & \cl{$\log n$} & \cl{$\log n$} & \cf{$1$} & \cs{$n$}\,$^*$ & \cf{$n$} \\
hash table & \cf{$1$}\,$^\dagger$ & \cf{$1$}\,$^\dagger$ & \cf{$1$}\,$^\dagger$ & \cs{$n$} & \cs{$n$} & \cl{$n \log n$} \\
binary heap & \cl{$\log n$} & \cx & \cx & \cf{$1$} & \cx & \cx \\
\bottomrule
\end{tabular}
\par\smallskip
{\footnotesize $^*$ $\log n$ with the policy-based \texttt{ordered\_set} \quad $^\dagger$ expected \quad --- not supported}
\caption{The cost of each operation, in $\mathcal{O}$-notation, for five ways to store a dynamic set (green: constant,
orange: logarithmic or $n \log n$, red: linear). No structure is best at everything; the balanced tree is the only one with
no red cell among the everyday operations, which is why \texttt{set} is the default for ``sorted and changing''.}
\label{fig:ds-costs}
\end{figure}

The table explains several habits of experienced competitors. When the data does not change after it is read, a
\emph{sorted array} beats every tree: sort once, then binary search. When you only ever need the smallest element, a
\emph{heap} beats a \cc{set} by a large constant factor. When order does not matter at all, a \emph{hash table} is fastest.
And when you need several operations that no single structure provides, you combine structures, as the rest of this chapter
shows repeatedly.

% ======================================================================================================================
\section[Stacks]{Stacks \hfill \bronze}\label{sec:stacks}

A \defn{stack} is a sequence where you can only add and remove at one end, the \emph{top}: the last element pushed is the
first popped (LIFO, last in, first out). Stacks model anything with \emph{nesting}: brackets, function calls, undo history,
the path of a depth-first search.

\begin{problem}{Balanced brackets}
Given a string of \texttt{()[]\{\}}, decide whether every bracket is closed by a matching bracket of the same type in the
right order.
\solution Scan the string. Push every opening bracket; for a closing bracket, the top of the stack must be the matching
opening bracket, which you pop. The string is balanced if no check fails and the stack is empty at the end
(Figure~\ref{fig:brackets}). Each character is pushed and popped at most once: $\Oh{n}$.
\end{problem}

\begin{figure}[htbp]
\centering
\begin{tikzpicture}[st/.style={cellc, minimum width=6mm, minimum height=5mm, font=\ttfamily\small}]
  % brackets are encoded as letters (p = paren, s = square, c = curly; capitals close them) because \foreach cannot hold
  % a bare "(" or "[" in its list
  \def\brp{(}\def\brP{)}\def\brs{[}\def\brS{]}\def\brc{\{}\def\brC{\}}
  \foreach \ch/\stack [count=\k] in {p/{p}, s/{p,s}, S/{p}, c/{p,c}, p/{p,c,p}, P/{p,c}, C/{p}, P/{}} {
    \pgfmathsetmacro{\x}{(\k-1)*1.35}
    \node[font=\ttfamily\normalsize\bfseries, text=fgblue] at (\x, 0.55) {\csname br\ch\endcsname};
    \draw[line width=0.6pt, fgink!60] (\x-0.4, -2.1) -- (\x+0.4, -2.1);
    \if\relax\detokenize\expandafter{\stack}\relax\else
      \foreach \s [count=\h] in \stack \node[st, fill=fgfill] at (\x, {-2.1+0.25+(\h-1)*0.5}) {\csname br\s\endcsname};
    \fi
  }
  \node[lbl, anchor=east] at (-0.5, 0.55) {read};
  \node[lbl, anchor=east] at (-0.5, -1.2) {stack after};
  \node[font=\sffamily\scriptsize\bfseries, text=fggreen!80!black, anchor=north] at (9.45, -2.2) {empty: balanced};
\end{tikzpicture}
\caption{Matching \texttt{([]\{()\})} with a stack. Each column shows the stack after reading the character above it; a
closing bracket always meets its partner on top, and the stack ends empty.}
\label{fig:brackets}
\end{figure}

The same scan evaluates arithmetic expressions (a stack of numbers and a stack of operators), checks that tags in an HTML-like
format are properly nested, and computes the depth of the deepest nesting. Recursion itself is a stack: every function call
pushes a frame and every return pops one, which is why a very deep recursion can run out of stack memory
(Section~\ref{sec:memory}) and why an explicit stack is the standard fix.

\subsection{The monotonic stack}

A stack whose contents are kept sorted, by popping everything that would break the order before each push, answers
``nearest'' questions in linear time. It is one of the most useful patterns in this book.

\begin{problem}{Nearest Smaller Values (CSES)}
For every position of an array, find the nearest position to its left with a smaller value (print its 1-indexed position,
or 0 if there is none). Limits: $n \le 2 \cdot 10^5$.
\end{problem}

Scan from left to right and keep a stack of candidate positions. When $a_i$ arrives, every candidate whose value is
$\ge a_i$ can never be the answer for any later position: $a_i$ is both closer and at least as small. So pop them all; what
remains on top is the answer for $i$; then push $i$. The stack's values increase from bottom to top at all times
(Figure~\ref{fig:monostack}).

\begin{figure}[htbp]
\centering
\begin{tikzpicture}
  % bars of a = 2 5 1 4 8 3 2 5 at x = i; height 0.32 cm per unit
  \foreach \v [count=\i from 0] in {2,5,1,4,8,3,2,5} {
    \fill[fgblue!25, draw=fgblue, line width=0.5pt] ({\i*1.1}, 0) rectangle ++(0.7, {\v*0.32});
    \node[font=\sffamily\small] at ({\i*1.1+0.35}, {\v*0.32+0.22}) {\v};
    \node[lbl] at ({\i*1.1+0.35}, -0.25) {\the\numexpr\i+1\relax};
  }
  % arcs below the axis, from each element to its nearest smaller element on the left
  \foreach \i/\j in {1/0, 3/2, 4/3, 5/2, 6/2, 7/6} {
    \pgfmathsetmacro{\depth}{0.25 + 0.17*(\i-\j)}
    \draw[->, line width=0.8pt, fgverm] ({\i*1.1+0.35}, -0.48) .. controls ({\i*1.1+0.35}, {-0.48-\depth})
      and ({\j*1.1+0.35}, {-0.48-\depth}) .. ({\j*1.1+0.35}, -0.52);
  }
  % answers and stack contents
  \node[lbl, anchor=east] at (-0.15, -1.55) {answer};
  \node[lbl, anchor=east] at (-0.15, -2.05) {stack (values)};
  \foreach \ans/\st [count=\i from 0] in {0/{2}, 1/{2 5}, 0/{1}, 3/{1 4}, 4/{1 4 8}, 3/{1 3}, 3/{1 2}, 7/{1 2 5}} {
    \node[font=\sffamily\small\bfseries, text=fgverm] at ({\i*1.1+0.35}, -1.55) {\ans};
    \node[font=\ttfamily\scriptsize] at ({\i*1.1+0.35}, -2.05) {\st};
  }
\end{tikzpicture}
\caption{Nearest smaller values of $2, 5, 1, 4, 8, 3, 2, 5$ (positions under the bars). Each arrow points to the answer.
The stack, bottom to top, only ever holds increasing values: when 3 arrives, 8 and 4 are popped for good, because 3 hides
them from everything to its right.}
\label{fig:monostack}
\end{figure}

\begin{code}
void solve() {
    int n;
    cin >> n;
    vi a(n), ans(n);
    cin >> a;
    vi st;  // indices; their values increase from the bottom to the top
    rep(i, n) {
        while (!st.empty() && a[st.back()] >= a[i])
            st.pop_back();  // can never be an answer again
        ans[i] = st.empty() ? 0 : st.back() + 1;  // 1-indexed, 0 if none
        st.push_back(i);
    }
    print(ans);
}
\end{code}

The inner \cc{while} can pop many elements at once, but every index is pushed once and popped at most once, so the whole
scan is $\Oh{n}$ (the potential argument of Section~\ref{sec:amortized}). Running the same scan from the right gives the
nearest smaller element on the right; flipping the comparison gives nearest larger. Together they give each element's
``reign'', the maximal window where it is the minimum, which is exactly what the contribution technique of
Section~\ref{sec:contribution} needed. The largest rectangle in a histogram, the number of visible buildings, and the
next warmer day are all one monotonic stack each.

\usetemplate{Monotonic stack (next greater / previous smaller)}{Ctrl+Alt+T, 7}{cpmonostack}

% ======================================================================================================================
\section[Queues and deques]{Queues and deques \hfill \silver}\label{sec:queues}

A \defn{queue} removes elements in the order they arrived (FIFO, first in, first out). Queues model waiting lines and,
above all, \emph{breadth-first processing}: BFS visits vertices in the order they were discovered, which is why it finds
shortest paths in unweighted graphs (Chapter~\ref{ch:graphs}). A \defn{deque} (double-ended queue) allows pushing and popping
at both ends in $\Oh{1}$.

\subsection{The sliding window maximum}

\begin{problem}{Window maximum}
For every window of $k$ consecutive elements, print its maximum ($n \le 2\cdot10^5$).
\end{problem}

Recomputing each window costs $\Oh{nk}$. The deque version keeps the indices of the elements that could still become a
window maximum, with their values decreasing from front to back. When $a_i$ arrives, every element at the back that is
$\le a_i$ can never be a maximum again (it is older and not larger), so pop it; then push $i$. If the front has slid out of
the window, pop it from the front. The front is now the maximum of the window (Table~\ref{tab:deque-trace}). Like the
monotonic stack, every index enters and leaves once: $\Oh{n}$ in total.

\begin{table}[htbp]
\centering
\caption{The monotonic deque on $1, 3, -1, -3, 5, 3, 6, 7$ with $k = 3$. Values are shown instead of indices.}
\label{tab:deque-trace}
\small
\begin{tabular}{c r l l r}
\toprule
\thd{$i$} & \thd{$a_i$} & \thd{what happens} & \thd{deque (front to back)} & \thd{window max} \\
\midrule
0 & 1 & push & 1 & \\
1 & 3 & pop 1 (smaller), push & 3 & \\
2 & $-1$ & push & 3, $-1$ & 3 \\
3 & $-3$ & push & 3, $-1$, $-3$ & 3 \\
4 & 5 & pop $-3$, $-1$, 3, push & 5 & 5 \\
5 & 3 & push & 5, 3 & 5 \\
6 & 6 & pop 3, 5, push & 6 & 6 \\
7 & 7 & pop 6, push & 7 & 7 \\
\bottomrule
\end{tabular}
\end{table}

\begin{code}
void solve() {
    int n, k;
    cin >> n >> k;
    vi a(n);
    cin >> a;
    deque<int> dq;  // indices in the window; values decrease front to back
    vi out;
    rep(i, n) {
        while (!dq.empty() && a[dq.back()] <= a[i])
            dq.pop_back();  // dominated by a[i]
        dq.push_back(i);
        if (dq.front() <= i - k) dq.pop_front();  // slid out of the window
        if (i >= k - 1) out.push_back(a[dq.front()]);
    }
    print(out);
}
\end{code}

\begin{recognize}
``For every window of length $k$'' or ``for every $i$, the best $j$ in $[i - k, i)$'': a monotonic deque. It also speeds
up dynamic programming with transitions from a sliding range, $dp_i = \min_{i-k \le j < i} dp_j + c_i$, from $\Oh{nk}$ to
$\Oh{n}$ (Chapter~\ref{ch:dpopt}).
\end{recognize}

\usetemplate{Sliding window min / max (monotonic deque)}{Ctrl+Alt+T, 6}{cpwindow}

% ======================================================================================================================
\section[Heaps]{Heaps \hfill \silver}\label{sec:heaps}

A \defn{heap} always knows its smallest (or largest) element. Inserting and removing the top cost $\Oh{\log n}$ each, and
nothing else is supported. That narrow interface is exactly what many algorithms need: Dijkstra's algorithm, Prim's
algorithm, event simulations and scheduling all repeat ``take the best remaining item, then add a few new items''.

\subsection{How a heap works}

A binary heap is a complete binary tree stored in an array (Section~\ref{sec:containers}) where every parent is at least as
large as its children (for a max-heap). To insert, put the new element in the first free slot at the bottom and let it
\emph{sift up}: while it is larger than its parent, swap them (Figure~\ref{fig:siftup}). To remove the top, move the last
element to the root and let it \emph{sift down}: while it is smaller than its larger child, swap them. Both walk one path
of the tree, so both cost $\Oh{\log n}$.

\begin{figure}[htbp]
\centering
\newcommand{\heapfig}[2]{% values of indices 0..8, index to highlight
  \begin{tikzpicture}[hv/.style={vtxs, minimum size=5.4mm}]
    \foreach \v [count=\i from 0] in {#1} {
      \pgfmathsetmacro{\lev}{int(floor(log2(\i+1)+0.001))}
      \pgfmathsetmacro{\pos}{int(\i+1-2^\lev)}
      \pgfmathsetmacro{\wid}{4.2/2^\lev}
      \pgfmathsetmacro{\xx}{(\pos+0.5)*\wid-2.1}
      \edef\st{\ifnum\i=#2 hot\else n\fi}
      \node[hv, \st] (n\i) at (\xx, -\lev*0.72) {\v};
    }
    \foreach \c/\p in {1/0, 2/0, 3/1, 4/1, 5/2, 6/2, 7/3, 8/3} \draw[ed] (n\p) -- (n\c);
  \end{tikzpicture}}
\begin{filmstrip}[2]
  \begin{stepbox}{1. put 10 in the first free slot}\heapfig{9,7,8,3,6,2,5,1,10}{8}\end{stepbox}
  \begin{stepbox}{2. $10 > 3$: swap with its parent}\heapfig{9,7,8,10,6,2,5,1,3}{3}\end{stepbox}
  \begin{stepbox}{3. $10 > 7$: swap again}\heapfig{9,10,8,7,6,2,5,1,3}{1}\end{stepbox}
  \begin{stepbox}{4. $10 > 9$: swap; it is the new root}\heapfig{10,9,8,7,6,2,5,1,3}{0}\end{stepbox}
\end{filmstrip}
\caption{Inserting 10 into a max-heap. The new element rises along one path from the bottom to the root; everything off that
path is untouched, so the cost is the height of the tree, $\Oh{\log n}$.}
\label{fig:siftup}
\end{figure}

\begin{fact}[Building a heap is linear]
Inserting $n$ elements one by one costs $\Oh{n \log n}$, but \cc{make_heap} (or constructing a \cc{priority_queue} from a
whole vector) builds a heap in $\Oh{n}$ by sifting down from the bottom level up: half the nodes are leaves and never move,
a quarter move at most one level, an eighth at most two, and $\sum_h h \cdot n / 2^{h+1} = \Oh{n}$.
\end{fact}

\subsection{Scheduling with a heap}

\begin{problem}{Room Allocation (CSES)}
A hotel has $n$ guests; guest $i$ stays from day $a_i$ to day $b_i$. Two guests can share a room if one leaves strictly
before the other arrives. Find the minimum number of rooms and a valid assignment. Limits: $n \le 2 \cdot 10^5$.
\solution Process guests by arrival. Keep a min-heap of the occupied rooms keyed by their guest's departure day. For a new
guest, look at the room that frees up earliest: if its guest left before today, reuse it; otherwise open a new room. Reusing
the earliest-freed room is never worse (an exchange argument), and a new room is opened only when every existing room is
occupied today, so the number of rooms equals the maximum number of guests present at once, which is a lower bound.
\end{problem}

\begin{code}
void solve() {
    int n;
    cin >> n;
    vector<array<ll, 3>> c(n);  // (arrival, departure, index)
    rep(i, n) cin >> c[i][0] >> c[i][1], c[i][2] = i;
    sort(all(c));
    minpq<pll> busy;  // (departure day, room) of the occupied rooms
    vi room(n);
    int rooms = 0;
    for (auto [a, b, i] : c) {
        if (!busy.empty() && busy.top().first < a) {  // earliest leaver gone
            room[i] = busy.top().second;
            busy.pop();
        } else room[i] = ++rooms;
        busy.push({b, room[i]});
    }
    print(rooms);
    print(room);
}
\end{code}

\subsection{The running median with two heaps}

\begin{problem}{Running median}
Numbers arrive one at a time. After each one, print the median of the numbers so far (the lower median when the count is
even).
\solution Keep the smaller half in a max-heap \texttt{low} and the larger half in a min-heap \texttt{high}, with
\texttt{low} holding the same number of elements or one more. The median is then the top of \texttt{low}
(Figure~\ref{fig:twoheaps}). Insert a new number into the half it belongs to, then move one top element across if the sizes
drifted apart. Every step costs $\Oh{\log n}$.
\end{problem}

\begin{figure}[htbp]
\centering
\begin{tikzpicture}
  \foreach \v [count=\i from 0] in {1,3,5,7,8,9,10,15} {
    \edef\st{\ifnum\i<4 cool\else warn\fi}
    \node[vtx, \st] (v\i) at (\i*1.05, 0) {\v};
  }
  \draw[line width=1.4pt, fgverm, dashed] (3.675, -0.75) -- (3.675, 0.75);
  \draw[brace] ($(v0.north west)+(0,0.12)$) -- ($(v3.north east)+(0,0.12)$) node[midway, above=7pt, font=\sffamily\scriptsize,
    text=fgblue] {\texttt{low}: max-heap, top 7};
  \draw[brace] ($(v4.north west)+(0,0.12)$) -- ($(v7.north east)+(0,0.12)$) node[midway, above=7pt, font=\sffamily\scriptsize,
    text=fgorange!80!black] {\texttt{high}: min-heap, top 8};
  \node[font=\sffamily\scriptsize\bfseries, text=fgverm] at (3.675, -1.05) {median = top of \texttt{low} = 7};
  \node[notel, anchor=west, text width=4.3cm] at (8.6, 0) {After the stream $5, 15, 1, 3, 8, 7, 9, 10$ the printed medians
    were $5, 5, 5, 3, 5, 5, 7, 7$.};
\end{tikzpicture}
\caption{Two heaps split the numbers at the median. Only the two tops are ever inspected, and they sit on either side of the
split, so the median is always at hand.}
\label{fig:twoheaps}
\end{figure}

\begin{code}
void solve() {
    int n;
    cin >> n;
    priority_queue<ll> low;  // the smaller half, largest on top
    minpq<ll> high;          // the larger half, smallest on top
    vi med;
    rep(i, n) {
        ll x;
        cin >> x;
        if (low.empty() || x <= low.top()) low.push(x);
        else high.push(x);
        // keep the sizes equal, or low one larger
        if (sz(low) > sz(high) + 1) high.push(low.top()), low.pop();
        else if (sz(high) > sz(low)) low.push(high.top()), high.pop();
        med.push_back(low.top());  // the lower median
    }
    print(med);
}
\end{code}

\begin{trick}[Deleting from a heap anyway]
When elements must also leave (a sliding median, where the oldest element expires), keep a second heap of ``deleted but
not yet removed'' values. Whenever the tops of the real heap and the deletion heap are equal, pop both. The heap then behaves
as if it supported deletion, at the same asymptotic cost. For the sliding median, a \cc{multiset} with an iterator to the
median is often simpler.
\end{trick}

% ======================================================================================================================
\section[Ordered sets and interval sets]{Ordered sets and interval sets \hfill \silver}\label{sec:ordered-sets}

A \cc{set} is the structure for questions about \emph{order} in a changing collection: the smallest element, the successor
of $x$, the element closest to $x$. It shines when the values are huge (coordinates up to $10^9$) and the collection
changes.

\begin{problem}{Closest occupied seat}
People take and leave seats in a row of length $10^9$. When a person arrives at seat $x$, report the distance to the nearest
occupied seat.
\solution Keep the occupied seats in a \cc{set}. For an arrival at $x$, \cc{s.lower_bound(x)} is the nearest occupied seat
at or to the right of $x$, and its predecessor (if any) is the nearest on the left. Each operation is $\Oh{\log n}$, and the
row's length never matters.
\end{problem}

\subsection{Sets of intervals}

Many problems maintain a collection of disjoint intervals: painted segments of a line, free blocks of memory, booked time
slots. Store them in a \cc{map} from left end to right end. To insert $[l, r]$, find the first interval that could overlap it
with \cc{lower_bound}, merge every overlapping interval into $[l, r]$ while erasing it, then insert the union
(Figure~\ref{fig:intervalset}). Every interval is inserted once and erased at most once, so $n$ insertions cost
$\Oh{n \log n}$ in total, however much merging happens.

\begin{figure}[htbp]
\centering
\begin{tikzpicture}[x=0.72cm]
  \numline{0}{17}{0}{}
  \node[lbl, anchor=east] at (-0.5, 1.75) {before};
  \ivl{1}{3}{1.75}{fgblue!50}{} \ivl{5}{7}{1.75}{fgverm!50}{} \ivl{9}{12}{1.75}{fgverm!50}{} \ivl{15}{16}{1.75}{fgblue!50}{}
  \node[lbl, anchor=east] at (-0.5, 1.15) {insert};
  \ivl{6}{10}{1.15}{fggreen!60}{}
  \node[lbl, anchor=east] at (-0.5, 0.55) {after};
  \ivl{1}{3}{0.55}{fgblue!50}{} \ivl{5}{12}{0.55}{fggreen!60}{} \ivl{15}{16}{0.55}{fgblue!50}{}
\end{tikzpicture}
\caption{Inserting $[6, 10]$ into a set of disjoint intervals. It overlaps $[5, 7]$ and $[9, 12]$ (red), which are erased and
replaced, together with it, by their union $[5, 12]$. The overlapping intervals are consecutive in the \texttt{map}, so one
\texttt{lower\_bound} and a short walk find them all.}
\label{fig:intervalset}
\end{figure}

The same structure, with each interval carrying a value, handles ``assign value $v$ to the range $[l, r]$'' operations: split
the intervals at $l$ and $r$, erase everything in between, insert one new interval. When the data is random this runs fast
and is known as the \emph{Chtholly tree}; with adversarial data, prefer a segment tree with lazy assignment
(Chapter~\ref{ch:segtree}).

% ======================================================================================================================
\section[Hash maps for counting]{Hash maps for counting \hfill \silver}\label{sec:hashmaps}

A hash table answers ``have I seen this key, and how often?'' in expected $\Oh{1}$. It is the natural structure for
frequency counting, for detecting duplicates, and for the classic ``two sum'' pattern: to find two elements with sum $x$, scan
the array and, for each $a_i$, ask whether $x - a_i$ was seen before. The prefix-sum version of the same idea counts
subarrays with a given sum (Section~\ref{sec:pigeonhole} and Chapter~\ref{ch:prefix}).

Two warnings from Chapter~\ref{ch:cpp} apply. Hash tables have a large constant factor, so when the keys are small integers,
an array indexed by the key is faster still. And a predictable hash can be attacked; use the randomised \cc{custom_hash}
from the library whenever others can see your code.

% ======================================================================================================================
\section[Augmenting simple structures]{Augmenting simple structures \hfill \gold}\label{sec:augment}

A structure can often store a little extra information per element that makes a new query cheap. This \defn{augmentation}
is one of the most powerful ideas in data structure design: segment trees, order-statistic trees and link-cut trees are all
augmented versions of simpler structures.

\paragraph{A stack with its minimum.} Store with each element the minimum of the stack from the bottom up to it. The minimum
of the whole stack is then the value stored with the top, in $\Oh{1}$, and popping never invalidates the stored minimums
below, because they only depend on what lies beneath.

\paragraph{A queue made of two stacks.} A queue can be built from two stacks: push onto an \texttt{in} stack; pop from an
\texttt{out} stack, and when \texttt{out} is empty, pour all of \texttt{in} into it, which reverses the order so that the
oldest element ends on top (Figure~\ref{fig:twostacks}). Each element is moved once, so every operation is $\Oh{1}$
amortized. Combined with the stack-with-minimum, this gives a \emph{queue with minimum}: a sliding window minimum that
works even when the window does not move one step at a time.

\begin{figure}[htbp]
\centering
\begin{tikzpicture}[st/.style={cellc, minimum width=14mm, minimum height=5.5mm, font=\ttfamily\small}]
  \newcommand{\stk}[4]{% x, title, list of "value/min", colour
    \begin{scope}[xshift=#1 cm]
      \draw[line width=0.7pt, fgink!70] (-0.8, 1.8) -- (-0.8, 0) -- (0.8, 0) -- (0.8, 1.8);
      \if\relax\detokenize{#3}\relax\else  % an empty stack draws no cells
        \foreach \v/\m [count=\h] in {#3} \node[st, fill=#4] at (0, {0.3+(\h-1)*0.55}) {\v\,\textcolor{fggray}{\scriptsize(\m)}};
      \fi
      \node[font=\sffamily\scriptsize\bfseries] at (0, -0.35) {#2};
    \end{scope}}
  \node[font=\sffamily\small\bfseries, anchor=west] at (-1.2, 2.45) {after push 4, 2, 5};
  \stk{0}{in}{4/4, 2/2, 5/2}{fgfill}
  \stk{2.1}{out}{}{fgfill}
  \draw[->, line width=1pt, fgblue] (3.4, 0.9) -- (4.8, 0.9) node[midway, above, font=\sffamily\scriptsize, text=fgblue] {pop};
  \node[font=\sffamily\small\bfseries, anchor=west] at (5.2, 2.45) {pop: pour \texttt{in} into \texttt{out}, pop the top};
  \stk{6.3}{in}{}{fgfill}
  \stk{8.4}{out}{5/5, 2/2}{fggreen!15}
  \node[notel, anchor=west, text width=3.4cm] at (9.6, 0.9) {4 was removed (the oldest). The queue is now 2, 5; its minimum
    is the smaller of the two top labels: 2.};
\end{tikzpicture}
\caption{A queue with minimum built from two stacks. Each entry stores its value and, in grey, the minimum of its stack from
the bottom up to it. Pouring reverses the order, so the oldest element is popped first.}
\label{fig:twostacks}
\end{figure}

\begin{code}
void solve() {
    int q;
    cin >> q;
    vector<pll> in, out;  // stacks of (value, minimum of the stack up to here)
    auto push = [](vector<pll>& st, ll x) {
        st.push_back({x, st.empty() ? x : min(x, st.back().second)});
    };
    while (q--) {
        int type;
        cin >> type;
        if (type == 1) {  // push x at the back
            ll x;
            cin >> x;
            push(in, x);
        } else if (type == 2) {  // pop the front
            if (out.empty())  // pour "in" into "out": the oldest ends on top
                while (!in.empty()) push(out, in.back().first), in.pop_back();
            out.pop_back();
        } else {  // minimum of the whole queue
            ll m = LLONG_MAX;
            if (!in.empty()) chmin(m, in.back().second);
            if (!out.empty()) chmin(m, out.back().second);
            print(m);
        }
    }
}
\end{code}

The same trick works for any \emph{associative} operation, not just the minimum: gcd, bitwise AND, even matrix products or
function composition. It is the simplest way to answer ``aggregate of a sliding window'' when the operation has no inverse, so
prefix sums cannot help.

% ======================================================================================================================
\section[Linked lists in arrays]{Linked lists in arrays \hfill \silver}\label{sec:linked}

\cc{std::list} is rarely worth using: it is slow to traverse and gives no random access. But the \emph{idea} of a linked list,
two arrays \cc{nxt[i]} and \cc{prv[i]} holding each element's neighbours, is the standard way to simulate deletions from a
sequence. Removing element $i$ is two assignments, \cc{nxt[prv[i]] = nxt[i]} and \cc{prv[nxt[i]] = prv[i]}, in $\Oh{1}$, and
the neighbours of any remaining element are available immediately.

\begin{problem}{Josephus Problem I (CSES)}
Children $1, \dots, n$ stand in a circle. Starting from child 1, every second child is removed until none remain. Print the
order of removal. Limits: $n \le 2\cdot10^5$.
\solution Keep a circular linked list in the array \cc{nxt}. The child after the current one is removed by linking the
current child to the next-but-one, then we skip ahead by one (Figure~\ref{fig:josephus}). Each removal is $\Oh{1}$, so the
whole simulation is $\Oh{n}$, where erasing from a \cc{vector} would cost $\Oh{n^2}$.
\end{problem}

\begin{figure}[htbp]
\centering
\begin{tikzpicture}
  \foreach \k in {1,...,7} {
    \pgfmathsetmacro{\ang}{90 - (\k-1)*360/7}
    \edef\st{\ifnum\k=2 dim\else\ifnum\k=4 dim\else\ifnum\k=6 dim\else n\fi\fi\fi}
    \node[vtx, \st] (c\k) at (\ang:1.6) {\k};
  }
  \draw[->, line width=0.9pt, fgblue] (c1) to[bend left=38] (c3);
  \draw[->, line width=0.9pt, fgblue] (c3) to[bend left=38] (c5);
  \draw[->, line width=0.9pt, fgblue] (c5) to[bend left=38] (c7);
  \draw[->, line width=0.9pt, fgblue] (c7) to[bend left=20] (c1);
  \foreach \k in {2, 4, 6} \draw[line width=1.2pt, fgverm] ($(c\k)+(-0.2,-0.2)$) -- ($(c\k)+(0.2,0.2)$)
    ($(c\k)+(-0.2,0.2)$) -- ($(c\k)+(0.2,-0.2)$);
  \node[notel, anchor=west, text width=5.8cm] at (2.6, 0) {After three removals (2, 4, 6, crossed out) the links skip
    them: \texttt{nxt[1] = 3}, \texttt{nxt[3] = 5}, \texttt{nxt[5] = 7}, \texttt{nxt[7] = 1}. The next child removed is
    \texttt{nxt[7] = 1}, and the full order for $n = 7$ is $2, 4, 6, 1, 5, 3, 7$.};
\end{tikzpicture}
\caption{The Josephus circle as a linked list in an array. Removing a child only redirects one link.}
\label{fig:josephus}
\end{figure}

\begin{code}
void solve() {
    int n;
    cin >> n;
    vi nxt(n + 1);
    rep(i, 1, n + 1) nxt[i] = i % n + 1;  // a circular linked list in an array
    vi order;
    int cur = 1;  // the child just before the next one to be removed
    rep(k, n) {
        int victim = nxt[cur];
        order.push_back(victim);
        nxt[cur] = nxt[victim];  // unlink the victim in O(1)
        cur = nxt[cur];          // skip one child
    }
    print(order);
}
\end{code}

Removing every $k$-th child instead of every second one needs ``find the $k$-th remaining child'', which a linked list
cannot do quickly; that is Josephus Problem II, solved with an order-statistic tree or a Fenwick tree in $\Oh{n \log n}$.

\begin{trick}[Deleting in reverse]
Many ``delete elements one by one and answer after each deletion'' problems become ``add elements one by one'' when processed
backwards, and adding is usually easier: a DSU can merge but not split, a sorted structure can gain neighbours cheaply, a
maximum can only grow (Section~\ref{sec:reverse}).
\end{trick}

\begin{summary}[Chapter summary]
\begin{itemize}
  \item List the operations and their frequencies first; the structure follows from them.
  \item Stacks handle nesting; the monotonic stack answers every ``nearest smaller / larger'' question in $\Oh{n}$.
  \item Queues give arrival order; the monotonic deque gives every sliding-window extreme in $\Oh{n}$.
  \item Heaps repeatedly give the best item; two heaps give a running median; lazy deletion adds removal.
  \item Ordered sets answer neighbour queries on huge coordinates; a \texttt{map} of intervals keeps disjoint ranges.
  \item Augment structures with extra information (a minimum per stack entry), and simulate deletions with
        \texttt{nxt}/\texttt{prv} arrays.
\end{itemize}
\end{summary}

\begin{practice}
\ppgroup{Stacks, queues and heaps}
\begin{pplist}
  \item Nearest Smaller Values
  \item Room Allocation
  \item Movie Festival II
  \item Towers
  \item Traffic Lights
  \item Maximum Subarray Sum II
\end{pplist}
\ppgroup{Ordered sets and simulation}
\begin{pplist}
  \item Concert Tickets
  \item Playlist
  \item Josephus Problem I
  \item Josephus Problem II
  \item Sliding Median
  \item Sliding Cost
\end{pplist}
\end{practice}
